\section{Mesoscopic consistency}\label{sec:meso}

The macroscopic proof uses the fact that the number of lines on the original board is of order $k$.
When few rows remain relative to $n$, the number of diagonals meeting the residual board
can be large relative to $k$. In this regime it is useful to consider how much mass they carry.
A diagonal from which many squares have been removed may have very small mass.

Section~\ref{sec:B5T01} imposes an upper bound on the mass of each diagonal. Under this hypothesis
and the quantitative conditions on rows and columns, the estimates are expressed in terms of
the masses and operator norms. Theorem~\ref{thm:meso} in
Section~\ref{sec:B5T05} thus gives the same consistency bound with constants
independent of $n/k$.
